%! Logarithmic Expressions — Unit 5, Lesson 5.1 — Exit Ticket
\documentclass[10pt]{article}
\usepackage{precalculus-article}
\usepackage{precalculus-boxes}
\newcommand{\TallMath}[1]{$\displaystyle #1\rule[-1.4em]{0pt}{3.2em}$}

\begin{document}

\pageheader{Unit 5, Lesson 5.1}{Exit Ticket}
\namedateperiod

\vspace{0.16in}

\begin{enumerate}[itemsep=20pt]

\item Convert to logarithmic form: $3^5 = 243$.

      Logarithmic form: \blank{5cm}

      \vspace{0.04in}
      Then evaluate: $\log_3 243 = $ \blank{2cm}

\item Evaluate without a calculator.

      $\log_{10} 10{,}000 = $ \blank{2.5cm} \qquad\qquad $\ln e^3 = $ \blank{2.5cm}

\item Between which two consecutive integers does $\log_2 50$ fall?
      Explain without a calculator.

      \vspace{0.04in}
      $2^{\blank{1cm}} = $ \blank{1.5cm} $< 50 <$ \blank{1.5cm} $= 2^{\blank{1cm}}$

      \vspace{0.04in}
      So $\log_2 50$ is between \blank{1cm} and \blank{1cm}.

      \vspace{0.04in}
      Explanation: \writeline

\end{enumerate}

\end{document}
